\section*{TD NSI :  Mission "L'Odyssée du Livreur"}

\textbf{Objectif de la mission :} \textit{Agent, votre mission est d'optimiser les livraisons de notre entreprise. Pour manipuler des coordonnées géographiques et des trajets fixes sans risque d'erreur, vous devez maîtriser une structure de données immuable : le \textbf{tuple}. Préparez votre itinéraire, car une erreur d'indexation et c'est la panne sèche !} 


\begin{Exercise}[title={Maîtrise des Tuples : Les fondamentaux }]
Répondez aux questions suivantes pour vérifier votre compréhension avant de passer à la suite de la mission.

\begin{enumerate}
    \item  Laquelle de ces syntaxes crée un \textbf{tuple} contenant uniquement le nombre 42 ?\\
    \hfill $\square$ \texttt{valeur = (42)} \hfill $\square$ \texttt{valeur = [42]} \hfill $\square$ \texttt{valeur = (42,)} \hfill

    
    \item  Soit le tuple \texttt{vitesse = (80, 110, 130)}. Quelle est la valeur de \texttt{vitesse[0]} ? \dotfill
    
    \item  Que renvoie l'instruction \texttt{vitesse[-1]} ? \dotfill
    
    \item  Quelle fonction permet d'obtenir le nombre d'éléments dans le tuple \texttt{vitesse} ? \dotfill
    
    \item  Si j'exécute \texttt{type(vitesse)}, quel résultat s'affiche dans la console ? \dotfill

    \item  On exécute : \texttt{couleurs = ("bleu", "vert")} puis \texttt{couleurs[1] = "rouge"}. 
    Citez le nom de l'erreur affichée par Python : \dotfill
    
    \item  Pourquoi cette erreur survient-elle ? (Complétez : le tuple est ............) \dotfill

    \item  Soit le tuple \texttt{config = (1920, 1080)}. Que vaut \texttt{len(config)} ? \dotfill
    
    \item  On souhaite parcourir le tuple suivant. Complétez la boucle pour afficher chaque ville :
            \begin{lstlisting}[style=pythonstyle, language=Python]
etapes = ("Tsimkoura", "Chirongui", "Poroani")
for .......................................................:
    print(ville)
\end{lstlisting}

    \item  On veut maintenant afficher l'index \textbf{et} la ville. Complétez :
            \begin{lstlisting}[style=pythonstyle, language=Python]
for i, v in ....................................................:
    print(f"Escale {i} : {v}")
\end{lstlisting}
    \item  Soit \texttt{data = (10, 20, 30)}. Que se passe-t-il si on demande \texttt{data[3]} ? \\
    .\dotfill

    \item  Comment tester si le nombre \texttt{20} est présent dans le tuple \texttt{data} ? \dotfill

    \item  On a \texttt{t1 = (1, 2)} et \texttt{t2 = (3, 4)}. Quel est le résultat de \texttt{t1 + t2} ? \dotfill

    \item  Soit \texttt{point = (5, 12)}. Si on écrit \texttt{x, y = point}, quelle est la valeur de \texttt{y} ? \dotfill

    \item  Observez ce code :
\begin{lstlisting}[style=pythonstyle, language=Python]
scores, total = (12, 15, 18), 0
for s in scores:
    total = total + s
\end{lstlisting}
    Quelle sera la valeur finale de \texttt{total} ? \dotfill

    \item  Dans le code précédent, quelle instruction permet d'accéder à la valeur \texttt{15} sans utiliser de boucle ? \\.\dotfill

    \item  Soit \texttt{mystere = (1, (2, 3), 4)}. Quel est le résultat de \texttt{len(mystere)} ? \dotfill

    \item  Quel symbole utilise-t-on pour délimiter un tuple ? \\
    \hfill $\square$ Des crochets \texttt{[]} \hfill $\square$ Des parenthèses \texttt{()} \hfill $\square$ Des accolades \texttt{\{\}} \hfill

    \item  Complétez ce code pour compter combien de fois le nombre \texttt{7} apparaît :
    \begin{lstlisting}[style=pythonstyle, language=Python]
valeurs, compteur = (7, 4, 7, 2), 0
for .............................................:
    if v == 7:
        compteur = ...........................................
\end{lstlisting}

    \item  Vrai ou Faux : On peut ajouter un élément à un tuple existant avec une méthode \texttt{.append()}. Justifiez.\\ 
    .\dotfill
\end{enumerate}
\end{Exercise}

\begin{Exercise}[title={Mission MayLivraison : Le défi du Drone }]
Le drone de l'entreprise \textit{MayLivraison} doit effectuer une tournée de livraison. Pour des raisons de sécurité, le drone doit impérativement \textbf{revenir à son point de départ} (cycle).

\begin{lstlisting}[style=pythonstyle, language=Python]
villes = {
    "Mamoudzou": (-12.781, 45.231), "Dapani": (-12.973, 45.174),
    "Sada": (-12.848, 45.104), "Chiconi": (-12.833, 45.118),
    "Dembeni": (-12.843, 45.187), "Bandrele": (-12.908, 45.192),
    "Chirongui": (-12.932, 45.148), "Boueni": (-12.910, 45.076),
    "Kani-Keli": (-12.956, 45.105), "Tsimkoura": (-12.912, 45.131),
    "Choungui": (-12.957, 45.129), "Poroani": (-12.891, 45.135),
    "Mbouini": (-12.983, 45.133), "Mzouazia": (-12.908, 45.089)
}
\end{lstlisting}

\begin{enumerate}
    \item  \textbf{Données :} Pourquoi est-il préférable d'utiliser des \textbf{tuples} pour stocker les coordonnées GPS plutôt que des listes ? \dotfill
    
    \item  \textbf{Accès :} Écrivez l'instruction permettant d'accéder à la \textbf{longitude} (2ème valeur) de "Sada" dans le dictionnaire \texttt{villes} : \dotfill
    
    \item  \textbf{Parcours :} On exécute le code : \texttt{for v in villes: print(v)}. Qu'affiche-t-il ?\\
    \hfill $\square$ Les couples (Nom, Coordonnées) \hfill $\square$ Uniquement les noms des villes \hfill $\square$ Uniquement les tuples de coordonnées
    \item  \textbf{Distance :} Complétez la fonction \textit{dist} calculant la distance euclidienne entre deux villes \textit{v1} et \textit{v2}:
    \begin{lstlisting}[style=pythonstyle, language=Python]
def dist(v1, v2, donnees):
    dx = .......................................................
    dy = .......................................................
    return (dx**2 + dy**2)**0.5\end{lstlisting}

    \item  \textbf{Le Mur de la Complexité :}
    \begin{enumerate}
        \item Combien de trajets différents existent pour 3 villes ? \dotfill ; pour 4 villes ? \dotfill ; pour 5 villes ? \dotfill
        \item Quelle opération mathématique (notée $n!$) permet de généraliser ce calcul pour $n$ villes ?\\ 
        .\dotfill
        \item Calculez le nombre de trajets pour 10 villes : \dotfill ; pour 20 villes : \dotfill
        \item Un processeur moderne teste $10^{9}$ trajets/sec. Combien de temps faut-il pour 50 villes ($3 \times 10^{64}$ trajets) ? \\
        .\dotfill\\
        .\dotfill\\
        \item \textbf{Conclusion :} Est-il faisable de trouver la solution parfaite par "force brute" (en testant toutes les solutions) ? Quel est l'intérêt d'une méthode gloutonne ? \\
        .\dotfill\\
        .\dotfill\\
    \end{enumerate}

    \item  \textbf{Analyse d'une instance (Cycle) :} 
    On considère le départ de \textbf{Chirongui} pour visiter : \textbf{Tsimkoura}, \textbf{Poroani}, \textbf{Mzouazia}.
    \begin{enumerate}
        \item Calculez la distance du trajet \texttt{Chi $\to$ Tsi $\to$ Por $\to$ Mzo $\to$ Chi} sachant que : \\
        $d(Chi, Tsi)=0.026$ ; $d(Tsi, Por)=0.021$ ; $d(Por, Mzo)=0.049$ ; $d(Mzo, Chi)=0.063$. \textbf{Total :} \dotfill
        \item \textbf{Solution Gloutonne :} En partant de \textbf{Chirongui}, quelle est la ville la plus proche entre Tsimkoura et Poroani ? \dotfill .\\
        Donnez le trajet glouton complet : \dotfill\\
        .\dotfill
        \item \textbf{Défi :} Trouvez un trajet (cycle) plus court que le glouton pour ces 4 villes : \dotfill
        .\dotfill\\
        .\dotfill\\
        \item \textbf{Conclusion :} La méthode gloutonne donne-t-elle toujours la meilleure solution possible ? \dotfill
    \end{enumerate}

    \item  \textbf{Algorithme :} Complétez le code pour trouver la ville la plus proche :
    \begin{lstlisting}[style=pythonstyle, language=Python]
def ville_proche(actuelle, liste_villes, coords):
    d_min = float("inf")
    prochaine = None
    for v in liste_villes:
        d = ..........(........., ......., .........)
        if d < .............:
            d_min = .............
            prochaine = .............
    return prochaine
    \end{lstlisting}

    \item  \textbf{Synthèse :} Codez la fonction \texttt{tournee\_gloutonne(depart, liste\_villes, coords)}.
    \textit{Aide (Pseudo-algo) :
    \begin{enumerate}
    \item Créer une liste \texttt{visitees} commençant par le départ.
    \item Créer une copie de \texttt{villes} nommée \texttt{a\_voir}.
    \item Tant qu'il reste des villes dans \texttt{a\_voir} :
    \begin{enumerate}
        \item Trouver la \texttt{prochaine} avec la fonction \texttt{ville\_proche}.
        \item L'ajouter aux \texttt{visitees} et la retirer de \texttt{a\_voir}.
    \end{enumerate}
    \item Ajouter le \texttt{depart} à la fin des \texttt{visitees} pour boucler le cycle.
    \end{enumerate}}
    
    \begin{lstlisting}[style=pythonstyle, language=Python]
def tournee_gloutonne(depart, villes, coords):

    .............................................................
    .............................................................
    .............................................................
    .............................................................
    .............................................................
    .............................................................
    .............................................................
    .............................................................
    .............................................................
    \end{lstlisting}
\end{enumerate}
\end{Exercise}


% \begin{Exercise}[title={Le Réseau Routier : Parcours Terrestre }]
% On utilise les distances routières (clés = tuples) :
% \begin{lstlisting}[style=pythonstyle, language=Python]
% cartes = {
%     ("Mtsamboro", "Dzoumogne"): 9.2, ("Dzoumogne", "Koungou"): 8.3,
%     ("Koungou", "Mamoudzou"): 7.5, ("Mamoudzou", "Dembeni"): 7.2,
%     ("Dembeni", "Bandrele"): 7.2, ("Sada", "Kani-Keli"): 12.2
% }
% \end{lstlisting}

% \begin{enumerate}
%     \item  Quel est le type de la clé \texttt{("Sada", "Kani-Keli")} ? \dotfill
%     \item  Écrivez une fonction \texttt{get\_dist(v1, v2)} qui renvoie la distance entre deux villes dans \texttt{dist\_routes}.
%     \item  Calculer le trajet \texttt{Mtsamboro $\to$ Dzoumogne $\to$ Koungou $\to$ Mamoudzou}. \\
%     \textbf{Distance :} \dotfill
%     \item  \textbf{Glouton :} De \texttt{Dzoumogne}, le drone a le choix entre \texttt{Mtsamboro} (9.2km) et \texttt{Koungou} (8.3km). Quel est le choix glouton ? \dotfill
%     \item  Proposez une amélioration du réseau (un nouveau tuple) pour raccourcir les trajets Nord-Sud : \dotfill
% \end{enumerate}
% \end{Exercise}

\begin{Comprendre}
 \textbf{BILAN}
\begin{itemize}
    \item \textbf{Tuple :} Immuable, peut servir de clé de dictionnaire.
    \item \textbf{Glouton :} Choix local optimal, rapide mais pas toujours parfait globalement.
    \item \textbf{Complexité :} La force brute (tout tester) échoue face à l'explosion factorielle ($n!$).
\end{itemize}
\end{Comprendre}